\documentclass[11pt,letterpaper,DIV=12]{scrartcl}

\usepackage[T1]{fontenc}
\usepackage[spanish,es-noshorthands]{babel}
\usepackage{lmodern}
\usepackage{graphicx}
\usepackage{booktabs}
\usepackage{tabularx}
\usepackage[hidelinks]{hyperref}
\usepackage{csquotes}
\usepackage[backend=biber,style=apa]{biblatex}

\addbibresource{references.bib}


\setcounter{secnumdepth}{0}

\begin{document}

\titlehead{\centering Instituto Tecnológico y de Estudios Superiores de Monterrey}
\subject{Inteligencia artificial avanzada para la ciencia de datos}
\title{Criterios de éxito}
\author{
  \begin{tabular}{l@{\hspace{2em}}l}
    Facundo Bautista Barbera & Emiliano Camacho Ponce \\
    A01066843 & A01712408 \\[1ex]
    Oswaldo Isaias Hernandez Santes & Iker Mejia Hernandez \\
    A01199004 & A01352358 \\[1ex]
    Alfredo Alejandro Soto Herrera & Jorge Manuel Oyoqui Aguilera \\
    A01711368 & A01711783
  \end{tabular}
}
\date{\today}

\maketitle


El éxito del proyecto se determina con indicadores del establo, comparados contra la línea base previa a la solución. La línea base se obtiene de los registros históricos de los 3 robots DeLaval VMS y de las evaluaciones visuales de BCS del personal. Los criterios técnicos del modelo se definen en 1.3.2.

\section{1. Criterios objetivos}

\begin{itemize}
  \item \textbf{Saturación de los robots:} las horas diarias con algún robot por encima del 90\% de utilización (línea base: 4 a 6 horas al día) se reducen al menos 35\%. Se mide con la utilización por robot y por hora en los registros del VMS. Evalúa: responsable del establo.
  \item \textbf{Espera en la manga:} el tiempo de espera promedio en horas pico (línea base: picos de 45 a 90 minutos) se mantiene en 20 minutos o menos. Se estima con el tiempo ocioso entre ordeños consecutivos, porque el VMS no registra la fila, y se valida con observación directa. Evalúa: responsable del establo y médico veterinario.
  \item \textbf{Flujo bimodal y producción:} las sesiones con flujo bimodal disminuyen 40\% y la producción diaria del hato no baja respecto a la línea base; la meta de referencia es recuperar de 150 a 200 kg de leche al día. Se mide con la producción y el perfil de flujo por ordeño registrados por el VMS. Evalúa: responsable del establo.
  \item \textbf{Vacas primíparas:} las vacas primíparas y de menor jerarquía mantienen en horas pico la frecuencia de ordeño definida por el establo. Se mide con la frecuencia de ordeño por vaca y por franja horaria. Evalúa: responsable del establo.
  \item \textbf{Evaluación del BCS:} el personal obtiene una calificación de BCS de cada vaca al menos cada dos semanas, sin manipular al animal y en menos tiempo que con la evaluación visual actual. Evalúa: médico veterinario y nutricionista.
  \item \textbf{Patologías metabólicas:} todas las vacas que entran al periodo seco con BCS $\geq$ 4.00 o que presentan BCS $\leq$ 2.25 se notifican a nutrición y veterinaria en la semana en que se detectan. Evalúa: médico veterinario y nutricionista.
  \item \textbf{Viabilidad reproductiva:} disminuye la proporción de vacas en posparto con BCS menor a 2.75. Es un indicador de largo plazo; su efecto en la tasa de preñez se observa después del cierre del proyecto. Evalúa: médico veterinario.
  \item \textbf{Confiabilidad de la información:} el personal atiende al menos 80\% de las alertas emitidas y no deja de usarlas por falsas alarmas. Evalúa: responsable del establo.
\end{itemize}

\section{2. Criterios subjetivos}

\begin{itemize}
  \item Las alertas de saturación llegan con anticipación suficiente y son claras para decidir cómo rebalancear el tráfico de vacas. Juicio: responsable del establo del CAETEC.
  \item La calificación automática de BCS es tan útil como la evaluación visual para decidir ajustes de nutrición y de manejo reproductivo. Juicio: médico veterinario del CAETEC.
  \item La solución puede operarse con el personal y el equipo actuales del establo (iPhone y registros del VMS) sin interrumpir la rutina de ordeño. Juicio: socio formador.
\end{itemize}

\section{3. Horizonte y criterio mínimo de éxito}

Los indicadores productivos y de salud necesitan semanas de operación para mostrar cambios. Por ello, al cierre del proyecto (diciembre de 2026) el éxito mínimo se verifica así:

\begin{itemize}
  \item Con datos históricos del CAETEC, demostrar que las alertas habrían anticipado los periodos de saturación registrados y estimar su efecto sobre los indicadores anteriores.
  \item Realizar al menos una prueba de captura con iPhone en el establo y obtener la validación del médico veterinario y del responsable del establo.
\end{itemize}

La verificación completa de las metas productivas se propone para un periodo piloto posterior de al menos 8 semanas, acordado con el CAETEC.

\printbibliography

\end{document}
